\documentclass{article}
\usepackage[T1]{fontenc}
\usepackage{lmodern}
\usepackage{graphicx}
\usepackage{amsmath,amssymb}
\usepackage[a4paper,margin=2cm]{geometry}
\usepackage[absolute,overlay]{textpos}
\setlength{\TPHorizModule}{1mm}
\setlength{\TPVertModule}{1mm}
\usepackage[indonesian]{babel}
\usepackage{hyperref}
\setlength{\emergencystretch}{2em}
% unit: Halaman 30

\begin{document}

% === Halaman 30 (llm) ===

4. \textbf{Cipherteks} (\textit{ciphertext}): pesan yang telah disandikan sehingga tidak bermakna lagi.

Tujuan: agar pesan tidak dapat dibaca oleh pihak yang tidak berhak.

Nama lain: \textbf{kriptogram} (\textit{cryptogram})

\begin{itemize}
    \item \textbf{Contoh:}
    \begin{itemize}
        \item \textbf{Plainteks:} \texttt{culik anak itu jam 11 siang}
        \item \textbf{Cipherteks:} \texttt{t\textasciicircum\$gfUi89rewoFpfdWqLMp[uTcxZ}
        \item \textbf{Kriptogram:} \texttt{t\textasciicircum\$gfU i89rewo FpfdWqLM p[uTcxZ}
    \end{itemize}
\end{itemize}

\end{document}
